
\begin{ledgroupsized}[r]{120mm}
\footnotesize
\pstart
\noindent\textbf{Überlieferung:}
\pend
\end{ledgroupsized}
\begin{ledgroupsized}[r]{114mm}
\footnotesize
\pstart
\parindent -6mm
\makebox[6mm][l]{\textit{L}}%
Auszüge mit Bemerkungen aus
\indexN{\textls{Giordani} (Giordano), Vitale 1633\textendash1711}\textsc{V.\;Giordani}, 
\cite{02110}\textit{De componendis gravium momentis dissertatio}, Rom~1687 und 
\cite{02109}\textit{Fundamentum doctrinae motus gravium}, 1.~Ausg. Rom~1688: 
LH~XXXV~14, 2~Bl.~89. 
Ein Blatt~2\textsuperscript{o};
italienisches Papier;
oberer und linker Rand beschnitten. 
Zwei Seiten.
\pend
\end{ledgroupsized}
\vspace{5mm}
\begin{ledgroup}
\footnotesize
\pstart
\noindent
\textbf{Datierungsgründe:}
Die vorliegenden Auszüge aus zwei Werken \indexN{\textls{Giordani} (Giordano), Vitale 1633\textendash1711}Vitale Giordanis lassen sich
in erster Näherung
aufgrund der Verwendung italienischen Papiers auf die Zeit des Italienaufenthaltes datieren: März 1689 bis März 1690.
Die Umstände von Leibnizens Begegnung mit \indexN{\textls{Giordani} (Giordano), Vitale 1633\textendash1711}Giordani und seinem Werk bieten zudem einen Terminus post quem für die Abfassung der Auszüge und erlauben die genauere Einkreisung ihrer Entstehung auf den Zeitraum Mitte November 1689 bis März 1690.
\pend
\pstart
Das erste persönliche Treffen zwischen Leibniz und \indexN{\textls{Giordani} (Giordano), Vitale 1633\textendash1711}Giordani fand wahrscheinlich Ende Oktober in Rom\indexO{Rom} statt: 
Siehe die Erläuterung zu Leibnizens erstem \cite{02119}Brief an Giordani (Ende Oktober bis Anfang November 1689, \textit{LSB} III,~4 N.\;216, bes. S.\;420).
Seinem darauffolgenden \cite{02120}Schreiben an Leibniz vom 11.~November 1689 legt Giordani ein Exemplar des \cite{02109}\textit{Fundamentum doctrinae motus gravium} von 1688 bei (N.\;217, bes. S.\;425\,f.).
Leibniz bedankt sich in seinem kurz vor der Abreise aus Rom\indexO{Rom} (am 21.\,(?) November 1689, siehe \cite{01236}\textit{Chronik}, S.\;98) 
verfassten \cite{02121}Brief 
bei \indexN{\textls{Giordani} (Giordano), Vitale 1633\textendash1711}Giordani
für sein \glqq neues Geschenk\grqq, das Leibniz auf der bevorstehenden Reise (\glqq in itinere\grqq) zu lesen gedenkt (N.\;218, hier S.\;427). 
Hiermit verstummt nach heutigem Kenntisstand der Briefwechsel.
Leibnizens Lektüre des \cite{02109}\textit{Fundamentum} kann daher frühestens Mitte November erfolgt sein; dieser Umstand liefert zugleich einen Terminus post quem für die Entstehung des Großteils der vorliegenden Auszüge N.\;\ref{58253}, der diesem Werk gewidmet ist.
\pend
\pstart
Was die \cite{02110}\textit{Dissertatio} von 1687 angeht, so gilt es anzumerken, dass 
\indexN{\textls{Giordani} (Giordano), Vitale 1633\textendash1711}Giordani
in seinem \cite{02120}Brief vom 11.~November Leibniz
\glqq\lbrack d\rbrack e honorifico Judicio tuo super mea \cite{02110}\textit{de momentis Dissertatione}\grqq\
dankt (N.\;217, S.\;423).
Diesem entweder mündlich oder
in der nicht erhaltenen Abfertigung des vorangeganenen Leibniz-Briefes (dessen Konzept schweigt N.\;216 über die \cite{02110}\textit{Dissertatio})
mitgeteilten Urteil muss eine, wenngleich kursorische, Lektüre des Opusculums durch Leibniz vorangegangen sein.
Dennoch erscheint die Annahme, Leibniz habe die ersten Absätze der Auszüge bereits nach dieser Lektüre, also vor dem Erhalt des \cite{02109}\textit{Fundamentum}, abgefasst,
in paläographischer Hinsicht als nicht plausibel: Weder der Schreibduktus noch die verwendete Tinte oder Feder lassen eine Unterbrechung im Textfluss erkennen.
Daraus ergibt sich für die Gesamtheit von N.\;\ref{58253} die vorgeschlagene Datierung: Mitte November 1689 bis März 1690.
\pend
\pstart
Leibniz exzerpiert aus dem ihm im November 1689 von \indexN{\textls{Giordani} (Giordano), Vitale 1633\textendash1711}Giordani geschenkten Exemplar der \cite{02109}ersten Ausgabe des \textit{Fundamentum} (Rom 1688).
In dieser Ausgabe fehlt die \glqq Beantwortung einiger Einwände, da sie noch nicht im Druck vorliegt\grqq, wie Giordani selbst an Leibniz schreibt (\textit{LSB} III,~4 \cite{02120}N.\;217, hier S.\;425). 
Gleichzeitig verspricht der Autor, nach Erscheinen des Werks Leibniz auch diesen Teil zukommen zu lassen, und zwar über \indexN{\textls{Ciampini}, Giovanni Giusto 1633\textendash1698}G.~Ciampini, der ihm bereits die Einwände mitgeteilt hatte (siehe \textsc{M.~Torrini}, \cite{02126}\textit{Dopo Galileo}, Florenz 1979, S.\;111).
Erst die zweite Ausgabe des Buchs (\cite{02118}\glqq Editio altera priori longe auctior, atque emendatior\grqq, Rom 1689) enthält genannte \glqq Responsiones\grqq\ (Anhang, S.\;82\textendash98).
Einen entsprechenden Auszug aus dieser \cite{02118}zweiten Ausgabe (S.\;81\textendash98) 
sowie deren Titelblatt und Widmung erhielt Leibniz vermutlich erst nach seiner Ankunft in Hannover\indexO{Hannover}
über \indexN{\textls{Ciampini}, Giovanni Giusto 1633\textendash1698}Ciampini und \indexN{\textls{Deseine}, François Jacques gest. 1715}François Deseine (wohl mit dessen \cite{02121}Brief aus Rom\indexO{Rom} vom 10.~Juni 1690, \textit{LSB} I,~5 N.\;339, hier S.\;586).
Tatsächlich ist das nachgesendete Material zusammen mit \indexN{\textls{Giordani} (Giordano), Vitale 1633\textendash1711}Giordanis \cite{02110}\textit{Dissertatio} und der Erstausgabe des \cite{02109}\textit{Fundamentum} zu Leibnizens Handexemplar gebunden (\textit{GWLB}, Nm-A 7049)
\pend
\pstart
Dass Leibniz aus der \cite{02109}\textit{ersten} Ausgabe exzerpiert hat, wird -- zusätzlich zu den Einzelheiten seiner Begegnung mit \indexN{\textls{Giordani} (Giordano), Vitale 1633\textendash1711}Giordani -- durch folgende Umstände bestätigt:
\indexN{\textls{Giordani} (Giordano), Vitale 1633\textendash1711}Giordani widmete 
die \cite{02109}Ausgabe von 1688 dem um 1689 verstorbenen 
\indexN{\textls{Howard of Worsop}, Thomas gest. 1689}Thomas Howard of Worsop 
(Leibniz hat die Stelle exzerpiert, siehe \refnote{FnC:35_14_02_089_Worsop}), während die \cite{02118}zweite Ausgabe
\indexN{\textls{Ottoboni}, Antonio 1646\textendash1720}Antonio Ottoboni zugeeignet ist.
Außerdem findet sich eine von Leibniz ausgezogene Stelle (siehe \refnote{FnC:35_14_02_089_finem}) nur in der \cite{02109}ersten Ausgabe des \textit{Fundamentum}, und zwar im Abschnitt \glqq Monita\grqq, bes. S.\;84\,f. im Schlussteil des Werks (Leibniz: \glqq sub finem operis\grqq); die Passage fehlt hingegen in der \cite{02118}zweiten Ausgabe.
\pend
\end{ledgroup}
\selectlanguage{latin}
\vspace{8mm}
\pstart
\noindent
\href{\Scan{LH_35_14_02_089r}}%
{\lbrack89~r\textsuperscript{o}\rbrack}\ 
\edtext{\indexN{\textls{Giordani} (Giordano), Vitale 1633\textendash1711}\textit{Vitalis Jordani Bitontini} 
\textit{in Romano Archigymnasio\indexS{Archigymnasium Romanum} necnon in 
regia Gallorum Academia\indexS{Academia regia Gallorum in Urbe} in \indexO{Rom}Urbe 
Matheseos professoris}
\cite{02110}\textit{De} 
\edtext{\textit{componendis gravium}}{%
\lemma{\textit{componendis}}%
\Bfootnote{%
\textit{(1)}~gravum
\textit{(2)}~\textit{gravium}~\textit{L}}}
\textit{momentis dissertatio.}
\textit{Romae typis Angeli Barnabò, 1687}\lbrack,\rbrack\ 
\edtext{folia 14.}{%
\lemma{folia}%
\Bfootnote{%
\textit{(1)}~85
\textit{(2)}~14
\textit{(3)}~14.~\textit{L}}}%
}{\lemma{\textit{Vitalis} \lbrack...\rbrack\ folia 14}%
\Cfootnote{\indexN{\textls{Giordani} (Giordano), Vitale 1633\textendash1711}\textsc{V.~Giordani}, \cite{02110}\textit{De componendis gravium momentis dissertatio}, Rom 1687, Titelseite.}}
\pend
\pstart
\edtext{In 
\edtext{praefatione: \textit{libellus superioribus mensibus editus}}{%
\lemma{praefatione:}%
\Bfootnote{%
\textit{(1)}~\textit{libellum superioribus mensibus editum}
\textit{(2)}~\textit{libellus superioribus mensibus editus}~\textit{L}}}
\textit{atque inscriptus} 
\cite{02000}\textit{Exegeses \greek{f}ysico-Mathematicae de momentis gravium, de vecte, ac de motu aequabiliter accelerato},
\indexN{\textls{Vanni}, Giovanni Francesco 1638\textendash1709}autor demonstrare aggressus est contra 
\indexN{\textls{Stevin} (Stevinus), Simon 1548\textendash1620}Stevinum et 
\indexN{\textls{Galilei} (Galilaeus, Galileus), Galileo 1564\textendash1642}Galilaeum
\textit{momentum partiale globi super plano declivi\indexS{momentum partiale globi super plano declivi} ad momentum suum totale\indexS{momentum totale globi}
non esse ut perpendiculum ad longitudinem plani declivis.}
\edtext{Cumque quidam putaret}{%
\lemma{Cumque}%
\Bfootnote{%
\textit{(1)}~Anonymus subtilis
\textit{(2)}~quidam putaret~\textit{L}}}
\textit{subtilem \indexN{\textls{Vanni}, Giovanni Francesco 1638\textendash1709}Anonymi doctrinam hoc} paralogismo 
\textit{laborare quod in ea componantur momenta gravium per additionem,\indexS{compositio momentorum gravium per additionem} cum sint revera}
\edtext{\textit{componend}\lbrack\textit{a}\rbrack}{%
\lemma{}%
\Bfootnote{%
componendae 
\textit{L ändert Hrsg. nach Vorlage}}}
\textit{per multiplicationem,\indexS{compositio momentorum gravium per multiplicationem}
dubitatum fuerit additio an Multiplicatio} praeferenda.}{%
\lemma{In praefatione \lbrack...\rbrack\ praeferenda}%
\applabel{FnC:35_14_02_089_Stevin}%
\Cfootnote{\cite{02110}a.\,a.\,O., S.~\lbrack III\rbrack, mit Auslassungen.
Siehe
\indexN{\textls{Vanni}, Giovanni Francesco 1638\textendash1709}\textsc{G.\,F.~Vanni} \lbrack anonym\rbrack, 
\cite{02000}\textit{Exegeses physico-mathematicae, de momentis gravium, de vecte, ac de motu aequabiliter accelerato}, Rom 1685,
sowie Leibnizens Rezension: N.\;\ref{41199} in diesem Band.
Zu Stevin und Galilei siehe
\indexN{\textls{Stevin} (Stevinus), Simon 1548\textendash1620}%
\textsc{S.~Stevin}, %
\cite{00196}\textit{De beghinselen der weeghconst}, Leiden 1586. %
Buch~1, Teil~2, Theorem~XI Prop.~XIX mit Koroll.~I, bes.\ S.\;40\textendash42
(\cite{02084}\textit{SPW} I, S.\;174\textendash179);
\indexN{\textls{Galilei} (Galilaeus, Galileus), Galileo 1564\textendash1642}%
\textsc{G.~Galilei}, %
\cite{00050}\textit{Discorsi}, %
Giornata Terza, De motu naturaliter accelerato; \glqq Aggiunta postuma dell'Autore\grqq\ %
nach dem Scholion zum 2.~Theorem. Der Zusatz erschien erst in der posthumen Werkausgabe: \cite{01084}\textit{Opere}, Bologna 1656, Bd.~II, S.\;132\lbrack ter\rbrack\textendash134 (\cite{00048}\textit{GO} VIII, S.\;216\textendash219 im Apparat).}}
\edtext{\indexN{\textls{Giordani} (Giordano), Vitale 1633\textendash1711}Autor 
\textit{a discipulis tota Europa dissitis}
\edtext{\textit{sententiam}}{\lemma{}\Bfootnote{\textit{sententiam} \textit{erg.~L}}}
\textit{rogatus}, hunc tractatum scripsit.}{%
\lemma{Autor \lbrack...\rbrack\ scripsit}%
\Cfootnote{\indexN{\textls{Giordani} (Giordano), Vitale 1633\textendash1711}\textsc{Giordani}, \cite{02110}\textit{De componendis gravium momentis}, S.~\lbrack III\rbrack f., mit Auslassungen.}}
\pend
\pstart
\edtext{%
\indexN{\textls{Casati}, Paolo 1617\textendash1707}\textit{Casatus} 
jam \textit{consideravit grave ex} duobus \textit{funiculis suspensum}\indexS{grave ex duobus funiculis suspensum}
\edtext{momentum\indexS{momentum gravis} exercere}{%
\lemma{momentum}%
\Bfootnote{%
\textit{(1)}~duplo
\textit{(2)}~exercere~\textit{L}}}
intelligi posse respectu unius solius quasi ex eo esset suspensum.}{%
\lemma{Casatus \lbrack...\rbrack\ suspensum}%
\Cfootnote{\cite{02110}a.\,a.\,O., S.\;1\,f. 
Siehe \indexN{\textls{Casati}, Paolo 1617\textendash1707}\textsc{P.~Casati}, 
\cite{02106}\textit{Mechanicorum libri octo}, Lyon 1684, lib.~I, cap.~XV.}}
\edtext{Et male ipsum inferri momentum totale\indexS{momentum totale gravis} \textit{non componi per additionem} ex duobus istis
\textit{sed temperari}\lbrack,\rbrack\ \textit{pag.~103.}\lbrack,\rbrack\
quin idem \textit{pag.~107 momenta} gravium\indexS{momentum gravis} \textit{per additionem componit.}\indexS{compositio momentorum gravium per additionem}}{%
\lemma{Et male \lbrack...\rbrack\ \textit{componit}}%
\Cfootnote{\indexN{\textls{Giordani} (Giordano), Vitale 1633\textendash1711}\textsc{Giordani}, \cite{02110}\textit{De componendis gravium momentis}, S.\;3 mit Änderungen.
Siehe \indexN{\textls{Casati}, Paolo 1617\textendash1707}\textsc{Casati}, \cite{02106}\textit{Mechanicorum libri octo}, lib.~I, cap.~XV, S.\;103 und S.\;107.}}
\edtext{Vult \textit{momenta ita temperari, ut}
\edtext{momentum}{\lemma{}\Bfootnote{momentum \textit{erg.~L}}}
\textit{diagonalis \textit{AN} repraesentet compositum\indexS{momentum compositum} ex \textit{AG} et \textit{AR} momentis.}}{%
\lemma{Vult \lbrack...\rbrack\ \textit{momentis}}%
\Cfootnote{\indexN{\textls{Giordani} (Giordano), Vitale 1633\textendash1711}\textsc{Giordani}, \cite{02110}\textit{De componendis gravium momentis}, S.\;4 mit Änderungen.}}
\pend
\begin{edfigure}%
\includegraphics[scale=1]{%
gesamttex/edit_VIII-4/images/LH_35_14_02_089_d1_089r.pdf}
\edcaption{\edtext{\lbrack\textit{Fig.~1}\rbrack}{%
\lemma{\lbrack\textit{Fig.~1}\rbrack}%
\Cfootnote[nonum]{Tlw. Übernahme der Vorlage: \cite{02110}a.\,a.\,O., S.\;4. Einen von Leibniz doppelt vergebenen Buchstaben \textit{A} ändert Hrsg. nach der Vorlage.}}}
\end{edfigure}
\pstart
\edtext{\textit{\textls{Definitio}. Gravia quorum aequales moles aequales habent gravitates} 
\edtext{seu \textit{pondera totalia}}{\lemma{}\Bfootnote{seu \textit{pondera totalia} \textit{erg.~L}}}
\textit{dicuntur \textls{ejusdem speciei}.}}{%
\lemma{\textit{\textls{Definitio}} \lbrack...\rbrack\ \textit{\textls{speciei}}}%
\Cfootnote{\cite{02110}a.\,a.\,O., S.\;5 mit Auslassungen.}}
\edtext{\textit{Gravia ejusdem speciei in eodem plano inclinato\indexS{planum inclinatum} inter se sunt ut momenta\indexS{momenta quae gravia exercent in plano} 
quae exercent in plano,} est prop.\;1.\lbrack,\rbrack\ 
quam demonstrat ex schol.
\edtext{\lbrack2.\rbrack}{\lemma{}\Bfootnote{1. \textit{L ändert Hrsg.}}}%
\;prop. de motu gravium 
\indexN{\textls{Galilei} (Galilaeus, Galileus), Galileo 1564\textendash1642}Galilaei,
et ex 
\indexN{\textls{Ghetaldi} (Getaldus), Marino 1566\textendash1626}Marini Ghetaldi prop.~2. et 3. \cite{00312}\textit{Archimedis promoti},}{%
\lemma{\textit{Gravia} \lbrack...\rbrack\ \textit{promoti}}%
\applabel{FnC:35_14_02_089_Ghetaldi}%
\Cfootnote{\cite{02110}a.\,a.\,O., S.\;6, theor.~I prop.~I.
Zu \indexN{\textls{Galilei} (Galilaeus, Galileus), Galileo 1564\textendash1642}%
Galilei siehe \cite{00050}\textit{Discorsi}, %
Giornata Terza; \glqq Aggiunta postuma dell'Autore\grqq\ %
nach dem Scholion zum 2.~Theorem (vgl. die Erläuterung zu \refnote{FnC:35_14_02_089_Stevin}). Leibnizens inkorrekter Verweis auf das \glqq schol. 1.~prop.\grqq\ geht auf Giordani zurück.
Zu Ghetaldi siehe 
\indexN{\textls{Ghetaldi} (Getaldus), Marino 1566\textendash1626}\textsc{M.~Ghetaldi},
\cite{00312}\textit{Promotus Archimedis}, Rom 1603, S.\;2\textendash4.}}
cum tamen res velut ex lemmis pateat, tantumque illud deest, 
quod non dedit definitionem momenti.\indexS{definitio momenti}
\edtext{Assumit igitur quod dixerat 
\indexN{\textls{Galilei} (Galilaeus, Galileus), Galileo 1564\textendash1642}Galilaeus\lbrack,\rbrack\
\textit{pondus totale gravis\indexS{pondus totale gravis} ad momentum\indexS{momentum gravis in plano inclinato} in plano} inclinato \textit{esse 
ut longitudo\indexS{longitudo plani inclinati}} plani inclinati \textit{ad perpendiculum}.}{%
\lemma{Assumit \lbrack...\rbrack\ \textit{perpendiculum}}%
\Cfootnote{\indexN{\textls{Giordani} (Giordano), Vitale 1633\textendash1711}\textsc{Giordani}, \cite{02110}\textit{De componendis gravium momentis}, S.\;6 mit Änderungen. 
Siehe \indexN{\textls{Galilei} (Galilaeus, Galileus), Galileo 1564\textendash1642}\textsc{Galilei}, \glqq Aggiunta postuma dell'Autore\grqq.}}
\pend
\pstart
\edtext{\textit{Theorema 4tum} est: 
\textit{absurdum est momentum\bracks{,} aequale composito ex momentis duorum gravium\indexS{momentum compositum} in} eodem \textit{plano inclinato}\bracks{,}
\textit{oriri ex continua multiplicatione} eorundem 
\edtext{\lbrack\textit{momentorum}\rbrack.}{%
\lemma{}\Bfootnote{gravium \textit{L ändert Hrsg. nach Vorlage}}}%
}{%
\lemma{\textit{Theorema} \lbrack...\rbrack\ \lbrack\textit{momentorum}\rbrack}%
\Cfootnote{\indexN{\textls{Giordani} (Giordano), Vitale 1633\textendash1711}\textsc{Giordani}, \cite{02110}\textit{De componendis gravium momentis}, S.\;11 mit Änderungen.}}
Sed quis hoc somniavit?
\pend
\pstart
Obiter
\edtext{in appendice asserit demonstrationem compositionis rationum\indexS{compositio rationum} non esse facilem
nec ab aliis praestitum ante Editionem \cite{02116}\textit{Euclidis} sui 
\edlabel{35_14_02_089_2a}%
\edtext{}{%
\xxref{35_14_02_089_2a}{35_14_02_089_2b}%
\lemma{\textit{restituti}. \lbrack/\rbrack\ \textit{Fundamentum}}%
\Cfootnote{Die Absätze sind durch eine waagerechte Linie getrennt.}}%
\textit{restituti}.}{%
\lemma{in appendice \lbrack...\rbrack\ restituti}%
\applabel{FnC:35_14_02_089_Euclide}%
\Cfootnote{\cite{02110}a.\,a.\,O., S.\;13.
Siehe \indexN{\textls{Giordani} (Giordano), Vitale 1633\textendash1711}\textsc{Giordani}, \cite{02116}\textit{Euclide restituto}, Rom~1680 (2.~Ausg. 1686), \glqq Annotazione\grqq\ und Beweise nach lib.~VI, prop.~XVIII, S.\;261\textendash271. Leibnizens Exemplar der zweiten Ausgabe des \textit{Euclide}
ist im Besitz der GWLB Hannover (Leibn.~Marg.~115).}}
\pend
\pstart
\edtext{%
\edtext{%
\cite{02109}\textit{Fundamentum\edlabel{35_14_02_089_2b} doctrinae motus gravium}
\textit{et comparatio momentorum gravis in planis sejunctis ad gravitationes}
\textit{quibus} corpus \textit{plana concurrentia premit}.
\textit{Geometrice restituta a \indexN{\textls{Giordani} (Giordano), Vitale 1633\textendash1711}Vitali Jordano Bitontino} \textit{Rom.~1688}.
fol. \textit{ex typographia Johannis Jacobi Komarek Bohemi apud Angelum Custodem}.%
}{\lemma{\textit{Fundamentum} \lbrack...\rbrack\ \textit{Custodem}}%
\Cfootnote{%
\indexN{\textls{Giordani} (Giordano), Vitale 1633\textendash1711}\textsc{Giordani},
\cite{02109}\textit{Fundamentum}, Titelseite mit Änderungen.}}%
}{\lemma{}%
\Afootnote{%
\textit{Am Rand:} %
\indexN{\textls{Giordani} (Giordano), Vitale 1633\textendash1711}Autor citat suum \cite{02106}\textit{Euclidem restitutum} et asserit in eo 
a se primum demonstrata\lbrack m\rbrack\textsuperscript{\lbrack a\rbrack} 
\indexS{doctrina proportionum}proportionum doctrinam.\textsuperscript{\lbrack b\rbrack}
Idem sub finem hujus operis promittit opus absolutum de motu gravium 
una cum sua mechanica.\textsuperscript{\lbrack c\rbrack} 
\\[1.5ex] %
{\footnotesize 
\textsuperscript{\lbrack a\rbrack}~demonstrata \textit{L ändert Hrsg.} \quad 
\textsuperscript{\lbrack b\rbrack}~asserit \lbrack...\rbrack\ doctrinam:
\indexN{\textls{Giordani} (Giordano), Vitale 1633\textendash1711}\textsc{Giordani},
\cite{02109}\textit{Fundamentum doctrinae motus gravium}, 1.~Ausg. Rom~1688, S.\;22.
Zu \indexN{\textls{Giordani} (Giordano), Vitale 1633\textendash1711}Giordanis
\cite{02116}\textit{Euclide restituto} siehe die Erläuterung zu \refnote{FnC:35_14_02_089_Euclide}.\quad 
\textsuperscript{\lbrack c\rbrack}~sub finem \lbrack...\rbrack\ mechanica: \indexN{\textls{Giordani} (Giordano), Vitale 1633\textendash1711}\textsc{Giordani}, \textit{Fundamentum}, monitum~IV, S.\;85.}}}
\pend
\pstart
Aggreditur probare et defendere 
\edtext{principium \indexN{\textls{Stevin} (Stevinus), Simon 1548\textendash1620}Stevini et 
\indexN{\textls{Galilei} (Galilaeus, Galileus), Galileo 1564\textendash1642}Galilaei.}{%
\lemma{principium Stevini et Galilaei}%
\Cfootnote{Siehe die Erläuterung zu \refnote{FnC:35_14_02_089_Stevin}.}}
\edtext{Librum 
\edtext{dedicat \textit{Domino}}{%
\lemma{dedicat}%
\Bfootnote{%
\textit{(1)}~\textit{Thomae}
\textit{(2)}~\textit{Domino}~\textit{L}}}
\indexN{\textls{Howard of Worsop}, Thomas gest. 1689}\textit{Thomae Howard de Worsop}\lbrack,\rbrack\
\indexN{\textls{Howard}, Henry VI.~duke of Norfolk 1677\textendash1684}\textit{ducis Norfolciae secundogenito} 
qui Romae\indexO{Rom} ab ipso ante 15 annos Mathesin audierat.}{%
\lemma{Librum \lbrack...\rbrack\ audierat}%
\applabel{FnC:35_14_02_089_Worsop}%
\Cfootnote{\cite{02109}a.\,a.\,O., S.~\lbrack III\rbrack\textendash\lbrack VI\rbrack.}}
\pend
\begin{edfigure}%
\includegraphics[scale=1]{%
gesamttex/edit_VIII-4/images/LH_35_14_02_089_d2_089r.pdf}
\edcaption{\edtext{\lbrack\textit{Fig.~2}\rbrack}{%
\lemma{\lbrack\textit{Fig.~2}\rbrack}%
\Cfootnote[nonum]{\cite{02109}a.\,a.\,O., S.\;1. Den Buchstaben \textit{G} ergänzt Hrsg.}}}
\end{edfigure}

\pstart
\edtext{%
\indexN{\textls{Torricelli} (Torricellius), Evangelista 1608\textendash1647}Torricellius 
ita rem probare instituit: 
\textit{Si grave\indexS{grave} circumferatur a} diametro \textit{\textit{FA} seu \textit{FH}
manente puncto \textit{F}}\lbrack,\rbrack\ \textit{tunc momentum totale\indexS{momentum totale gravis} ejus} 
seu \textit{momentum quod habet in \textit{H} ad momentum quod habet in \textit{A} 
ut \textit{FH} ad \textit{FL} hoc est ut \textit{DF} ad \textit{FI} vel \textit{BC} ad \textit{CE}.}
\indexN{\textls{Giordani} (Giordano), Vitale 1633\textendash1711}Autor obscuritatem hic invenit}{%
\lemma{Torricellius \lbrack...\rbrack\ invenit}%
\Cfootnote{\cite{02109}a.\,a.\,O., S.\;1 mit Änderungen;
siehe \indexN{\textls{Torricelli} (Torricellius), Evangelista 1608\textendash1647}%
\textsc{E.~Torricelli}, \textit{De motu gravium}, lib.~I, prop.~II, Lemma
(\cite{00304}\textit{Opera geometrica}, Florenz 1644, S.\;101\,f.;
\cite{00108}\textit{TO}~II, S.\;106f).}}
sed nulla est, certum enim est eodem modo moveri grave, in \textit{A}
\edtext{super plano inclinato,\indexS{planum inclinatum}}{\lemma{}\Bfootnote{super plano inclinato \textit{erg.~L}}}
ac si abesset planum inclinatum et grave conaretur in extremo semidiametro 
\edlabel{35_14_02_089_1a}%
\edtext{}{%
\xxref%
{35_14_02_089_1a}{35_14_02_089_1b}%
\lemma{\textit{FA}}%
\Bfootnote{%
\textit{(1)}~\textit{Gravitas}
\textit{(2)}~Autor~\textit{L}}}%
\textit{FA}.
\pend
\pstart
\indexN{\textls{Giordani} (Giordano), Vitale 1633\textendash1711}Autor\edlabel{35_14_02_089_1b}
hoc melius demonstrare aggreditur\lbrack:\rbrack\
\edtext{%
\edtext{\textit{Gravit}\lbrack\textit{as}\rbrack}{%
\lemma{}\Bfootnote{Gravitatis \textit{L ändert Hrsg. nach Vorlage}}}
\textit{totalis\indexS{gravitas totalis} est propensio} gravis \textit{si liberum relinqueretur.}}{%
\lemma{\textit{Gravit}\lbrack\textit{as}\rbrack}%
\Cfootnote{%
\indexN{\textls{Giordani} (Giordano), Vitale 1633\textendash1711}\textsc{Giordani}, \textit{Fundamentum}, def.~III, S.\;3.}}
\pend
\pstart
\edtext{\textit{Duo gravia contraponderantia}\indexS{gravia controponderantia} autor dicit \textit{librari, cum vero aequalia habent momenta 
\edlabel{35_14_02_089_3a}%
\edtext{}{%
\xxref{35_14_02_089_3a}{35_14_02_089_3b}%
\lemma{\textit{aequilibrari.}}%
\Bfootnote{%
\textit{(1)}~Ex
\textit{(2)}~Potentia~\textit{L}}}%
aequilibrari.}}{%
\lemma{Duo \lbrack...\rbrack\ aequilibrari}%
\Cfootnote{\cite{02109}a.\,a.\,O., def.~IV, monitum~I, S.\;5, mit Änderungen.}}
\pend
\begin{edfigure}%
\includegraphics[scale=1]{%
gesamttex/edit_VIII-4/images/LH_35_14_02_089_d3_089r.pdf}
\edcaption{\edtext{\lbrack\textit{Fig.~3}\rbrack}{%
\lemma{\lbrack\textit{Fig.~3}\rbrack}%
\Cfootnote[nonum]{\cite{02109}a.\,a.\,O., S.\;6.}}}
\end{edfigure}
\pstart
\edlabel{35_14_02_089_4a}%
\edtext{}{%
\xxref{35_14_02_089_4a}{35_14_02_089_4b}%
\lemma{Potentia \lbrack...\rbrack\ quartam}%
\Cfootnote{\cite{02109}a.\,a.\,O., exemplum~II, S.\;6.}}%
Potentia\edlabel{35_14_02_089_3b} sustinens\indexS{potentia sustinens corpus} 
\edtext{corpus\indexS{corpus ex vecte suspensum} \textit{A} ex vect\lbrack e\rbrack\ \textit{BD} suspensum ne descendat est}{%
\lemma{corpus \textit{A}}%
\Bfootnote{%
\textit{(1)}~ne descendat fit
\textit{(2)}~ex \bars{vectes \textit{ändert Hrsg.}} \textit{BD}
\textit{(a)}~ne
\textit{(b)}~suspensum ne descendat
\textit{(aa)}~opus
\textit{(bb)}~est~\textit{L}}}
aequalis momento
\edtext{peculiari}{\lemma{}\Bfootnote{peculiari \textit{erg.~L}}}
gravis ex \textit{CD} suspensi, altera vero pars sustinetur a fulcimento\indexS{fulcimentum} \textit{B}.
Ita autor pag.~6 in exemplo secundo post definitionem 
\edtext{quartam.\edlabel{35_14_02_089_4b} Hoc}{%
\lemma{quartam.}%
\Bfootnote{%
\textit{(1)}~(+~posset dicere, momentum peculiare absolutum. Sed in eo est error, \bars{quo \textit{streicht Hrsg.}} enim
\textit{(2)}~Hoc~\textit{L}}}
enim ejus sunt verba:
\edtext{\textit{Ex aliquo puncto D in vecte \textit{BC} sumto pendeat grave \textit{A} et 
unum vectis extremum incumbat fulcimento\indexS{fulcimentum} \textit{B}, alterum vero extremum \textit{C} 
ab aliqua potentia\indexS{potentia} sustineatur. Jam gravitatis pars ea quae sustinetur a potentia in \textit{C}, 
repraesentat momentum\indexS{momentum gravis} quod habet grave \textit{A} ex \textit{D} suspensum, 
quodque dicitur momentum peculiare gravis\indexS{momentum gravis} \textit{A} ex puncto \textit{D} suspensi.
Altera pars quae sustinetur a fulcimento\indexS{fulcimentum} \textit{B} 
dicitur momentum impeditum seu momentum resistentiae.\indexS{momentum resistentiae}}}{%
\lemma{\textit{Ex} \lbrack...\rbrack\ \textit{resistentiae}}%
\Cfootnote{\cite{02109}a.\,a.\,O., S.\;6.}}
Hoc ille\lbrack,\rbrack\ sed non video quod ita ratiocinari possit, cum alia atque alia sit 
\edtext{vis quo}{\lemma{}\Bfootnote{vis \bars{vis \textit{streicht Hrsg.}} quo~\textit{L}}}
\textit{C} 
\edtext{premitur quo propius}{%
\lemma{premitur quo}%
\Bfootnote{%
\textit{(1)}~gravius
\textit{(2)}~propius~\textit{L}}}
est ipsi \textit{D}\lbrack,\rbrack\ et cum ipsi \textit{D} supponitur est totalis.
Et revera tunc \textit{B} non sustinet.
\pend
\pstart
\href{\Scan{LH_35_14_02_089v}}%
{\lbrack89~v\textsuperscript{o}\rbrack}\ 
\edtext{Ipse mox fatetur \indexN{\textls{Giordani} (Giordano), Vitale 1633\textendash1711}autor si potentia sustineat 
\edtext{sub gravi nihil}{%
\lemma{sub}%
\Bfootnote{%
\textit{(1)}~pondere
\textit{(2)}~gravi
\textit{(a)}~non
\textit{(b)}~nihil~\textit{L}}}
sustinere fulcimentum.\indexS{fulcimentum}}{%
\lemma{Ipse \lbrack...\rbrack\ fulcimentum}%
\Cfootnote{\cite{02109}a.\,a.\,O., S.\;6.}}
Forte ergo voluit pressiones fulcimentorum 
esse reciproce ut distantias,\indexS{pressiones reciprocae ut distantias} sed male se 
\edtext{explicuit. Praeterea}{%
\lemma{explicuit.}%
\Bfootnote{%
\textit{(1)}~Nec hoc quicquam commune habet cum momento
\textit{(a)}~suspension
\textit{(b)}~suspensionis ex \textit{D}. Et certum est si abstrah
\textit{(2)}~Praeterea~\textit{L}}}
hoc debebat probare.
\pend
\begin{edfigure}%
\includegraphics[scale=1]{%
gesamttex/edit_VIII-4/images/LH_35_14_02_089_d4_089v.pdf}
\edcaption{\edtext{\lbrack\textit{Fig.~4}\rbrack}{%
\lemma{\lbrack\textit{Fig.~4}\rbrack}%
\Cfootnote[nonum]{\cite{02109}a.\,a.\,O., S.\;8, zweites Diagramm.}}}
\end{edfigure}
\pstart
\edtext{%
Momentum Gravis \textit{G} in \textit{E} positi sic ut pondera \textit{A} et \textit{G} 
sint in aequilibrio\indexS{aequilibrium} aequatur momento totali gravis\indexS{momentum totale gravis} \textit{A}, quia cum remoto pondere \textit{G}
totali suo momento\indexS{momentum quo corpus premit planum} premeret planum \textit{BC} nunc premit nullo modo.}{%
\lemma{Momentum \lbrack...\rbrack\ modo}%
\Cfootnote{\cite{02109}a.\,a.\,O., exemplum~IV, S.\;8.}}
Liceat \indexN{\textls{Giordani} (Giordano), Vitale 1633\textendash1711}autori nostro sic uti voce \textls{momenti}. Sed intelligendum est 
momentum relativum\indexS{momentum relativum} respectu scilicet ejus 
\edtext{cujus exercitium}{%
\lemma{cujus}%
\Bfootnote{%
\textit{(1)}~vim
\textit{(2)}~exercitium~\textit{L}}}
impediunt. Hinc infert tales conclusiones:
\edtext{%
\textit{Momentum minoris gravis\indexS{momentum gravis suspensi} ex majore distantia suspensi pondere totali ipsius minoris gravis est majus.
Momentum majoris gravis\indexS{momentum gravis suspensi} ex minore distantia suspensi ponderi totali ejusdem gravis est aequale. 
Momenti minoris gravis\indexS{momentum gravis suspensi} ex majore distantia suspensi}
\edtext{\textit{ponderi totali}}{%
\lemma{\textit{ponderi}}%
\Bfootnote{%
\textit{(1)}~ejusdem
\textit{(2)}~\textit{totali}~\textit{L}}}
\textit{majoris gravis est aequale.}}{%
\lemma{Momentum \lbrack...\rbrack\ aequale}%
\Cfootnote{\cite{02109}a.\,a.\,O., S.\;9.}}
\pend
\begin{edfigure}%
\includegraphics[scale=1]{%
gesamttex/edit_VIII-4/images/LH_35_14_02_089_d5_089v.pdf}
\edcaption[0.25em]{\edtext{\lbrack\textit{Fig.~5}\rbrack}{%
\lemma{\lbrack\textit{Fig.~5}\rbrack}%
\Cfootnote[nonum]{\cite{02109}a.\,a.\,O., S.\;14.}}}
\end{edfigure}

\pstart
\edtext{%
Si \textit{gravia \textit{A} et \textit{B}} inaequalia \textit{sint in aequilibrio\indexS{gravia inaequalia in aequilibrio} erit \textit{CE} ad \textit{ED}, 
ut grave \textit{B} ad grave \textit{A}}
ex 6. et 7.~prop. \indexN{\textls{Archimedes} 287\textendash212 v. Chr.}Archim. 
\cite{01010}\textit{De aequip.}
id \textit{est ut pondus totale gravis \textit{B} ad pondus totale gravis \textit{A}}
ex \textit{prop.~2. et 3. 
\indexN{\textls{Ghetaldi} (Getaldus), Marino 1566\textendash1626}Marini Ghetaldi in 
\cite{00312}\textit{Archimede promoto}.}
Ideo si \textit{CE} sit unitas 
\textit{productum ex multiplicatione distantiae \textit{ED}
in pondus totale\indexS{pondus totale}} 
\edtext{\textit{gravi}\lbrack\textit{s}\rbrack}{%
\lemma{}%
\Bfootnote{%
gravitatis 
\textit{L ändert Hrsg. nach Vorlage}}}
\edtext{\textit{B} dabit}{%
\lemma{\textit{B}}%
\Bfootnote{%
\textit{(1)}~seu
\textit{(2)}~sumta
\textit{(3)}~\bars{dabit \textit{streicht Hrsg.}}
\textit{(4)}~dabit~\textit{L}}}
\textit{pondus totale}
\edtext{\textit{gravi}\lbrack\textit{s}\rbrack}{%
\lemma{}%
\Bfootnote{gravitatis \textit{L ändert Hrsg. nach Vorlage}}}
\edlabel{35_14_02_089_5a}%
\edtext{}{%
\xxref{35_14_02_089_5a}{35_14_02_089_5b}%
\lemma{\textit{A}.}%
\Bfootnote{%
\textit{(1)}~Et
\textit{(2)}~(+~Omnia~\textit{L}}}%
\textit{A}.}{%
\lemma{Si gravia \lbrack...\rbrack\ \textit{A}}%
\Cfootnote{\cite{02109}a.\,a.\,O., S.\;14 mit Änderungen. 
Der Verweis auf \indexN{\textls{Archimedes} 287\textendash212 v. Chr.}\textsc{Archimedes}, 
\cite{01010}\textit{De planorum aequilibris libri duo}, lib.~I, prop.~VI\textendash VII fehlt in der Vorlage.
Zu Ghetaldi\indexN{\textls{Ghetaldi} (Getaldus), Marino 1566\textendash1626} siehe die Erläuterung zu \refnote{FnC:35_14_02_089_Ghetaldi}.}}
\pend
\begin{edfigure}[0.5em]%
\includegraphics[scale=1]{%
gesamttex/edit_VIII-4/images/LH_35_14_02_089_d6_089v.pdf}
\edcaption{\lbrack\textit{Fig.~6}\rbrack}
\end{edfigure}
\pstart
(+~Omnia\edlabel{35_14_02_089_5b} haec ita quidem dici possunt sed parvum usum habent.
\edtext{Nam si}{%
\lemma{Nam}%
\Bfootnote{%
\textit{(1)}~sive
\textit{(2)}~pondus
\textit{(3)}~si~\textit{L}}}
\textit{A} et \textit{B} ponantur 
\edtext{aequalia et}{%
\lemma{aequalia}%
\Bfootnote{%
\textit{(1)}~habent m
\textit{(2)}~et~\textit{L}}}
aequidistantia\indexS{gravia aequalia et aequidistantia} 
momenta eorum secundum ratiocinationem autoris sunt ipsa pondera totalia.
Sed hinc patet non habere fixam sedem momenta 
\edtext{totalia, sed}{%
\lemma{totalia,}%
\Bfootnote{%
\textit{(1)}~ut
\textit{(2)}~sed~\textit{L}}}
pendere ab opposito, ut proinde sensus sit pondus relativum\indexS{pondus relativum} esse majus absoluto,\indexS{pondus absolutum}
\edtext{cum majus corpus sustinet.}{%
\lemma{cum}%
\Bfootnote{%
\textit{(1)}~plus
\textit{(2)}~majus corpus
\textit{(a)}~elevat
\textit{(b)}~sustinet~\textit{L}}}
Quod satisfacit modo loquendi vulgari\lbrack,\rbrack\ sed si momenta semper aestimes
a conatu descendendi\indexS{conatus descendendi} seu quantitate velocitatis\indexS{velocitas} in molem\indexS{moles} ductae,\indexS{quantitas velocitatis in molem ductae} non est his opus.~+)
\pend
\begin{edfigure}[0.5em]%
\includegraphics[scale=1]{%
gesamttex/edit_VIII-4/images/LH_35_14_02_089_d7_089v.pdf}
\edcaption[0.em]{\edtext{\lbrack\textit{Fig.~7}\rbrack}{%
\lemma{\lbrack\textit{Fig.~7}\rbrack}%
\Cfootnote[nonum]{\indexN{\textls{Giordani} (Giordano), Vitale 1633\textendash1711}\textsc{Giordani}, \cite{02109}\textit{Fundamentum}, S.\;25, erstes Diagramm.}}}
\end{edfigure}
\pstart
\edtext{%
Sit vectis\indexS{vectis} \textit{BDA} ex quo suspensum pondus \textit{C} in \textit{D}, et potentia\indexS{potentia sustinens pondus}
\edtext{in \textit{A}}{\lemma{}\Bfootnote{in \textit{A} \textit{erg.~L}}}
sustinens pondus
\edtext{\textit{C}\lbrack,\rbrack\ erit}{%
\lemma{\textit{C}}%
\Bfootnote{%
\textit{(1)}~in \textit{A}
\textit{(2)}~quae
\textit{(3)}~erit~\textit{L}}}
ea potentia tanta ut sustineat pondus\indexS{potentia sustinens pondus} \textit{F}, quod 
\edtext{sit ad \textit{C} ut \textit{DB} ad \textit{AB}.}{%
\lemma{sit ad}%
\Bfootnote{%
\textit{(1)}~\textit{E}, ut
\textit{(2)}~\textit{C} ut
\textit{(a)}~\textit{BD} ad \textit{AB}
\textit{(b)}~\textit{DB} ad \textit{AB}.~\textit{L}}}
Hoc 
\edtext{probat sic\lbrack:\rbrack}{%
\lemma{probat}%
\Bfootnote{%
\textit{(1)}~quia
\textit{(2)}~sic~\textit{L}}}
sumatur $E=C$, 
\edtext{erit mom.}{%
\lemma{erit}%
\Bfootnote{%
\textit{(1)}~mom. \textit{E} ex \textit{A}, ad mom.
\textit{(2)}~mom.~\textit{L}}}
\textit{C} ex \textit{D} ad mom. \textit{E} ex \textit{A}, ut \textit{BD}
\edtext{ad \textit{BA} ex propositione~6 
\indexN{\textls{Del Monte}, Guidobaldo 1545\textendash1607}Guidubaldi de Monte in 
\cite{01009}tractatu \textit{de libra}, ut jam}{%
\lemma{ad \textit{BA}}%
\Bfootnote{%
\textit{(1)}~jam \textit{F} ad \textit{E}
\textit{(2)}~\bars{ex \lbrack...\rbrack\ \textit{libra} \textit{erg.}}
ut jam~\textit{L}}}
momenta ex iisdem punctis suspensa simul ut Gravia,
ergo mom. \textit{F} ex \textit{A} est ad mom. \textit{E} ex \textit{A}
ut \textit{BD} ad \textit{BA}. Ergo mom. \textit{C} ex \textit{D} et mom. \textit{F} ex \textit{A} sunt aequales.
Si ergo potentia \textit{A} sustinere potest pondus \textit{C} ex \textit{D} 
seu est aequalis momento \textit{C} ex \textit{D}, erit et aequalis momento \textit{F} ex \textit{A}.}{%
\lemma{Sit vectis \lbrack...\rbrack\ ex \textit{A}}%
\Cfootnote{\cite{02109}a.\,a.\,O., theor.~II prop.~II, S.\;25.
\indexN{\textls{Del Monte}, Guidobaldo 1545\textendash1607}\textsc{G. Del~Monte}, 
\cite{01009}\textit{Mechanicorum liber}, Pesaro 1577, \textit{De libra}, prop.~6, Bl.~40~v\textsuperscript{o}.}}
\pend
\begin{edfigure}[1em][1em]    	%
\includegraphics[scale=1]{%
gesamttex/edit_VIII-4/images/LH_35_14_02_089_d8_089v.pdf}
\edcaption{\lbrack\textit{Fig.~8}\rbrack}
\end{edfigure}
\pstart
\edtext{%
\textit{Momenta gravium ejusdem speciei a vectium sejunctorum extremitatibus pendentium}
vult \textit{esse ut eorum pondera totalia\indexS{pondus totale}}
\edtext{seu momentum}{%
\lemma{seu}%
\Bfootnote{%
\textit{(1)}~pondus totale \textit{C} ad
\textit{(2)}~momentum~\textit{L}}}
\textit{C} ad momentum\indexS{momentum} \textit{N} ut pondus\indexS{pondus} \textit{C} ad pondus \textit{N}.
Considerat scilicet quod aequalibus potentiis suppositis\indexS{gravia aequalibus potentiis suppositis sustentata} directe, sint sustinenda.}{%
\lemma{\textit{Momenta} \lbrack...\rbrack\ sustinenda}%
\Cfootnote{\cite{02109}a.\,a.\,O., theor.~IV prop.~IV, S.\;28\,f.}}
\edtext{Hinc enim ostendit momenta duorum gravium\indexS{momenta duorum gravium}
ex extremitatibus duorum vectium appensorum non esse 
\edtext{statim}{\lemma{}\Bfootnote{statim \textit{erg.~L}}}
ut extremitates ex quibus pendent nisi eadem sit gravium proportio\indexS{proportio gravium} et distantia.\indexS{distantia gravium}}{%
\lemma{Hinc \lbrack...\rbrack\ distantia}%
\Cfootnote{\cite{02109}a.\,a.\,O., coroll.~I, S.\;30.}}
\pend
\pstart
\edtext{Demum demonstrat 
\indexN{\textls{Giordani} (Giordano), Vitale 1633\textendash1711}autor 
\textit{momentum liberum\indexS{momentum liberum gravis} gravis in plano inclinato\indexS{planum inclinatum} aequale esse momento peculiari ejusdem gravis, 
quod solum tangens planum illud inclinatum a vecte vel libra totaliter sustinetur.}}{%
\lemma{Demum \lbrack...\rbrack\ \textit{sustinetur}}%
\Cfootnote{\cite{02109}a.\,a.\,O., theor.~VI prop.~VI, S.\;37.}}
Atque ita 
\edtext{%
\edtext{\indexN{\textls{Torricelli} (Torricellius), Evangelista 1608\textendash1647}Torricellii demonstratione}{%
\lemma{Torricellii}%
\Bfootnote{%
\textit{(1)}~Theorema pe
\textit{(2)}~demonstratione~\textit{L}}}
perficit sed ego vim non intelligo\lbrack,\rbrack\ est 
\edtext{prop.~\lbrack7\rbrack}{%
\lemma{}%
\Bfootnote{6 \textit{L ändert Hrsg.}}}
\edtext{\lbrack pag.\rbrack~38.}{%
\lemma{}%
\Bfootnote{fol. \textit{L ändert Hrsg.}}}%
}{%
\lemma{Torricellii \lbrack...\rbrack\ \lbrack pag.\rbrack~38}%
\Cfootnote{\cite{02109}a.\,a.\,O., theor.~VII prop.~VII, S.\;38\,f.}}
Nec vacat meditari.
\pend
\begin{edfigure}%
\includegraphics[scale=1]{%
gesamttex/edit_VIII-4/images/LH_35_14_02_089_d9_089v.pdf}
\edcaption{%
\edtext{\lbrack\textit{Fig.~9}\rbrack}{%
\lemma{\lbrack\textit{Fig.~9}\rbrack}%
\Cfootnote[nonum]{Vorlage sind die Diagramme \cite{02109}a.\,a.\,O., S.\;54 und S.\;55 mit Änderungen.}%
\lemma{}%
\Afootnote[nonum]{%
\textit{Unterhalb des Diagramms \lbrack\textit{Fig.~9}\rbrack}:
Totum erroneum}}%
}
\end{edfigure}
\pstart
\edtext{%
\edtext{%
Prop.~XV. pag.~54. \textit{Si grave sphaericum\indexS{grave sphaericum duobus planis sustentatum} duobus planis declivibus\indexS{planum declive}} \textit{DC}, \textit{CB}
\textit{in quemcunque angulum \textit{DCB} coeuntibus sustineatur}\lbrack,\rbrack\
\textit{ex centro autem} \textit{A} \textit{et ex contactibus} cum rectis \textit{D} et \textit{B} 
\textit{cadant ad horizontem} \textit{EGCF} \textit{perpendiculares} \textit{DE}, \textit{AG}, \textit{BF}\lbrack,\rbrack\
\textit{gravitationes quibus grave plana illa premit sunt in reciproca ratione cum perpendicularibus
quae de contactibus}
\edtext{\textit{cadunt} (hoc est gravitatio}{%
\lemma{\textit{cadunt}}%
\Bfootnote{%
\textit{(1)}~\textit{, et pondus totale gravis,}
\textit{(2)}~(+
\textit{(3)}~(hoc est 
\textit{(a)}~pondus quo premitur
\textit{(b)}~gravitatio~\textit{L}}}
qua premitur planum\indexS{gravitatio qua premitur planum} \textit{DC} est ad gravitationem qua premitur planum \textit{BC},
ut \textit{BF} ad \textit{DE}) \textit{et pondus}
\edtext{\textit{totale gravis} \textit{Z}}{%
\lemma{\textit{totale}}%
\Bfootnote{%
\textit{(1)}~gravitatis \textit{Z}
\textit{(2)}~\textit{gravis} \textit{Z}~\textit{L}}}
\textit{ad unam dictarum}}{%
\lemma{}%
\Afootnote{%
\textit{Am Rand}: %
(+~Subtilitatem non satis attingit quam ego explicui\bracks{:}\textsuperscript{\bracks{a}} 
planum diruptum iri si \textit{B}\textsuperscript{\lbrack b\rbrack}
non sit firmius quam
ab \indexN{\textls{Giordani} (Giordano), Vitale 1633\textendash1711}autore assignatum,
exigua plani firmitas\indexS{firmitas plani} debet sustinere posse potentiam liberam\indexS{potentia libera gravis descendentis}
descendentis per \textit{DC}, quae sit ad totam ut \textit{DE} ad \textit{DC}
seu ut \textit{NE} ad \textit{EF}~\lbrack+\rbrack)
\\[1.5ex] %
{\footnotesize 
\textsuperscript{\lbrack a\rbrack} explicui: Siehe Leibnizens \textit{Demonstratio geometrica regulae apud Staticos receptae}, \cite{01023}\textit{Acta Eruditorum} 1685, S.\;501--505, bes. S.\,505; hier als N.\;\ref{49409}, \refpassage{AE_1685_501-505_18a}{AE_1685_501-505_18b}.
\textsuperscript{\lbrack b\rbrack} \textit{B} \textit{erg.~L}}}}
\edtext{\textit{gravitationum} ut}{%
\lemma{\textit{gravitationum}}%
\Bfootnote{%
\textit{(1)}~\textit{est ut ho}
\textit{(2)}~ut~\textit{L}}}
in \textit{D} vel in \textit{B} \textit{est ut horizontalis} \textit{EF} 
\textit{intercepta inter dictas perpendiculares} \textit{DE}, 
\edtext{\textit{BF} ut}{%
\lemma{\textit{BF}}%
\Bfootnote{%
\textit{(1)}~ut
\textit{(2)}~ad
\textit{(3)}~ut~\textit{L}}}
pars horizontalis \textit{intercepta inter perpendicularem ex centro et ex contactu alterne sumtam}
seu grave \textit{Z} est ad gravitationem\indexS{gravitatio gravis} in \textit{D} ut \textit{EF} ad \textit{FG},
ad gravitationem vero in \textit{B} ut \textit{EF} ad \textit{EG}.}{%
\lemma{Prop.~XV. \lbrack...\rbrack\ \textit{EG}}%
\Cfootnote{\cite{02109}a.\,a.\,O., theor.~XIII prop.~XV, S.\;54\,f. mit Änderungen.}}
Haec ita probare aggreditur satis ingeniose:
\edtext{%
fingit vectem\indexS{vectis}
\edtext{\textit{DB} et grave}{%
\lemma{\textit{DB} et}%
\Bfootnote{%
\textit{(1)}~mobile suspensum
\textit{(2)}~grave~\textit{L}}}
suspensum\indexS{grave suspensum} ex puncto \textit{H} (in quod cadit linea directionis)\lbrack,\rbrack\ supponit\lbrack,\rbrack\
\edtext{nec satis probat quantum nunc judicare possum\lbrack,\rbrack}{\lemma{}\Bfootnote{nec \lbrack...\rbrack\ possum \textit{erg.~L}}}
eodem modo tunc a vecte plana premi,\indexS{pressio vectis in planum} ut nunc a globo.\indexS{pressio globi in planum} Hinc
\edtext{jam (contracto}{%
\lemma{jam}%
\Bfootnote{%
\textit{(1)}~ita ratiocinatur:
\textit{(2)}~(contracto~\textit{L}}}
autoris argumento quod prolixe admodum 
\edtext{exponit) ex}{%
\lemma{exponit)}%
\Bfootnote{%
\textit{(1)}~pro
\textit{(2)}~ex
\textit{(3)}~ex~\textit{L}}}
3~coroll. ad 2.~prop. 
\indexN{\textls{Del Monte}, Guidobaldo 1545\textendash1607}Guidubaldi de Monte in \cite{01009}tr. 
\edtext{\textit{de Vecte}; assumit}{%
\lemma{\textit{de Vecte};}%
\Bfootnote{%
\textit{(1)}~infert
\textit{(2)}~assumit~\textit{L}}}
potentiam quae sustinet vectem\indexS{potentia quae sustinet vectem} in \textit{D} esse ad eam quae sustinet in \textit{B}
ut \textit{HB} ad \textit{HD} seu ut \textit{GF} 
\edtext{ad \textit{EG}.}{%
\lemma{}%
\Bfootnote{ad \bars{ad \textit{streicht Hrsg.}} \textit{EG}.~\textit{L}}}
Probat autem esse \textit{GF} ad \textit{GE} ut \textit{BF} ad \textit{DE} hoc modo:
anguli \textit{EDO} et \textit{ADC} sunt aequales\lbrack,\rbrack\ ergo demto communi \textit{ODC}
erit $ADO=EDC$. Ergo $DAO=DCE$. Ergo $ED:DC\squaredots$
\edtext{$OD:DA$ (ob}{%
\lemma{$OD:DA$}%
\Bfootnote{%
\textit{(1)}~. Et simi
\textit{(2)}~(ob~\textit{L}}}
angulum ad \textit{O} rectum). Et similiter $CB:BF$ ut \textit{AD} seu \textit{AB} ad \textit{BL}\lbrack,\rbrack\
est autem $DC=BC$ (ob contactum in \textit{D} et \textit{B}). Ergo fiet $ED: BF\squaredots OD\ \text{seu}\ EG:BL\ \text{seu}\ GF$.
Quia autem \textit{EG} et \textit{GF} simul componunt \textit{EF}, 
et sunt \textit{EG} et \textit{GF} ut pressiones in plana,\indexS{pressio vectis in planum} utique \textit{EF} totam gravitatem\indexS{gravitas totalis globi}
globi repraesentabit.}{%
\lemma{fingit \lbrack...\rbrack\ repraesentabit}%
\Cfootnote{\cite{02109}a.\,a.\,O., S.\;55\textendash57. Siehe 
\indexN{\textls{Del Monte}, Guidobaldo 1545\textendash1607}\textsc{Del Monte}, 
\cite{01009}\textit{Mechanicorum liber}, \textit{De vecte}, prop.~2, coroll.~III, Bl.~40~v\textsuperscript{o}.}}
Consentit prior meis 
\edtext{nempe planum}{%
\lemma{nempe}%
\Bfootnote{%
\textit{(1)}~grave sustinet vim quae est ad
\textit{(2)}~pondus
\textit{(3)}~planum~\textit{L}}}
\textit{D} sustinet vim qua globus descendet in plano\indexS{vis qua globus descendet in plano} \textit{B}, quae est ad totam 
ut \textit{BF} ad \textit{BC} seu \textit{DC}, et grave \textit{B} sustinet vim quae est ad totam
ut \textit{DE} ad \textit{DC} seu ad \textit{BC} ergo vis in \textit{B} ad vim in \textit{D}
ut \textit{BF} ad \textit{DE}, seu ut \textit{GF} ad \textit{GE}, vis autem tota ut \textit{EF}.
\pend
\pstart
\edtext{Notat \indexN{\textls{Giordani} (Giordano), Vitale 1633\textendash1711}autor sub finem \cite{02109}operis
\indexN{\textls{Galilei} (Galilaeus, Galileus), Galileo 1564\textendash1642}Galilaeum 
\cite{00050}dial.~2 \textit{mechanicorum}
simili artificio investigasse \textit{cum pars saxi a subjecto plano\indexS{planum subjectum}}\lbrack,\rbrack\
\textit{pars ab extremitate alicujus vectis\indexS{vectis} sustinetur}\lbrack,\rbrack\ 
\textit{quota pars ponderis a vecte et plano sustineatur.}}{%
\lemma{Notat \lbrack...\rbrack\ \textit{sustineatur}}%
\applabel{FnC:35_14_02_089_finem}%
\Cfootnote{%
\indexN{\textls{Giordani} (Giordano), Vitale 1633\textendash1711}\textsc{Giordani}, \cite{02109}\textit{Fundamentum}, probl.~III prop.~XXXI, annotatio, S.\;84\,f.
Siehe \indexN{\textls{Galilei} (Galilaeus, Galileus), Galileo 1564\textendash1642}\textsc{Galilei}, \cite{00050}\textit{Discorsi}, Giornata Seconda, S.\;113 (\cite{00048}\textit{GO} VIII, S.\;155\,f.).}}
\pend
